\subsection{Addition von Vektoren}
\label{subsec:vector_algebra_addition}

\begin{DefinitionBox}{Vektoraddition}{vector_algebra_addition}
  Die \emph{Vektoraddition} \EnglishTerm{vector addition} zweier Vektoren
  $\vec a,\vec b\in\R^n$ erfolgt komponentenweise:
  \[
    \vec a+\vec b
    =\begin{pmatrix}a_1\\a_2\\\vdots\\a_n\end{pmatrix}
     +\begin{pmatrix}b_1\\b_2\\\vdots\\b_n\end{pmatrix}
    =\begin{pmatrix}
      \Cloze{2.4cm}{a_1+b_1}\\[0.25em]
      \Cloze{2.4cm}{a_2+b_2}\\[0.25em]
      \vdots\\[0.25em]
      \Cloze{2.4cm}{a_n+b_n}
    \end{pmatrix}.
  \]
  Vektoren unterschiedlicher Dimension können nicht addiert werden.
\end{DefinitionBox}

\begin{KeyPointBox}{Geometrische Vektoraddition}{vector_algebra_geometric_addition}
  Um die Summe zweier Vektoren geometrisch zu konstruieren, zeichnet man
  $\vec a$ und $\vec b$ mit gemeinsamem Anfangspunkt und ergänzt sie zu einem
  Parallelogramm. Die vom gemeinsamen Anfangspunkt ausgehende Diagonale ist der
  Summenvektor $\vec a+\vec b$.
  \begin{center}
    \CachedTikzFigure{content/resources/chapters/vector_algebra/tikz/vector_addition.tex}
  \end{center}
\end{KeyPointBox}

\begin{ExampleBox}{Vektoren rechnerisch und geometrisch addieren}{vector_algebra_addition_example_2d}
  \begin{minipage}[c]{0.52\linewidth}
    Gegeben seien die beiden Vektoren
    \[
      \vec a=\begin{pmatrix}1\\2\end{pmatrix},
      \qquad
      \vec b=\begin{pmatrix}3\\1\end{pmatrix}.
    \]
    Die komponentenweise Addition liefert
    \[
      \vec a+\vec b=
      \begin{ClozeContinuation}[12]{2}
        \begin{pmatrix}1+3\\2+1\end{pmatrix}
        =\begin{pmatrix}4\\3\end{pmatrix}
      \end{ClozeContinuation}
    \]
  \end{minipage}\hfill
  \begin{minipage}[c]{0.44\linewidth}
    \centering
    \CachedTikzFigure{content/resources/chapters/vector_algebra/tikz/vector_addition_example_2d.tex}
  \end{minipage}
  \par\smallskip
  Rechnerisch und geometrisch erhalten wir somit dasselbe Ergebnis.
\end{ExampleBox}

Da Vektoren komponentenweise addiert werden, lassen sich die Rechengesetze
der Addition von Zahlen unmittelbar auf die Vektoraddition übertragen.

\begin{TheoremBox}{Rechenregeln der Vektoraddition}{vector_algebra_addition_laws}
  Bei der Addition beliebiger Vektoren $\vec a,\vec b,\vec c\in\R^n$ gelten die
  üblichen algebraischen Rechenregeln:
  \[
    \begin{aligned}
      \vec a+\vec b
        &=\Cloze{5cm}{\vec b+\vec a}
        &&\qquad\text{(Kommutativgesetz)},\\[0.5em]
      \vec a+(\vec b+\vec c)
        &=\Cloze{5cm}{(\vec a+\vec b)+\vec c}
        &&\qquad\text{(Assoziativgesetz)}
    \end{aligned}
  \]
\end{TheoremBox}

Daher können Vektoren beim Umstellen von Gleichungen hinsichtlich Addition
und Subtraktion genauso wie Zahlen behandelt werden.

\subsubsection{Subtraktion als Addition des Gegenvektors}
\label{subsubsec:vector_algebra_subtraction}

Für die Subtraktion führen wir keine neue Rechenart ein. Sie baut auf der
Vektoraddition und dem Gegenvektor auf.

\begin{DefinitionBox}{Vektorsubtraktion}{vector_algebra_subtraction}
  Die \emph{Vektorsubtraktion} \EnglishTerm{vector subtraction} ist als
  Addition von $\vec a$ mit dem Gegenvektor $-\vec b$ definiert:
  \[
    \vec a-\vec b
    :=\vec a+(-\vec b)
    =\begin{pmatrix}a_1\\a_2\\\vdots\\a_n\end{pmatrix}
     +\begin{pmatrix}-b_1\\-b_2\\\vdots\\-b_n\end{pmatrix}
    =\begin{pmatrix}
      \Cloze{2.4cm}{a_1-b_1}\\[0.25em]
      \Cloze{2.4cm}{a_2-b_2}\\[0.25em]
      \vdots\\[0.25em]
      \Cloze{2.4cm}{a_n-b_n}
    \end{pmatrix}.
  \]
\end{DefinitionBox}

\begin{KeyPointBox}{Geometrische Vektorsubtraktion}{vector_algebra_geometric_subtraction}
  Gemäß der Definition schreiben wir die Differenz als Summe mit dem
  Gegenvektor:
  \[
    \vec a-\vec b=\vec a+(-\vec b).
  \]
  Zur geometrischen Konstruktion zeichnen wir, wie bei einer Summe üblich,
  $\vec a$ und $-\vec b$ mit gemeinsamem Anfangspunkt zu einem Parallelogramm.
  Dessen Diagonale ist der Differenzvektor $\vec a-\vec b$.
  \begin{center}
    \CachedTikzFigure{content/resources/chapters/vector_algebra/tikz/vector_subtraction_as_addition.tex}
  \end{center}
  Zugleich ist $\vec a-\vec b$ der Verbindungsvektor von der Spitze von
  $\vec b$ zur Spitze von $\vec a$.
\end{KeyPointBox}

\begin{ExampleBox}{Vektoren subtrahieren}{vector_algebra_componentwise_subtraction}
  \begin{minipage}[c]{0.52\linewidth}
    Gegeben seien die beiden Vektoren
    \[
      \vec a=\begin{pmatrix}4\\3\end{pmatrix},
      \qquad
      \vec b=\begin{pmatrix}1\\2\end{pmatrix}.
    \]
    Die Subtraktion liefert
    \[
      \vec a-\vec b=
      \begin{ClozeContinuation}[12]{2}
        \vec a+(-\vec b)
        =\begin{pmatrix}4-1\\3-2\end{pmatrix}
        =\begin{pmatrix}3\\1\end{pmatrix}
      \end{ClozeContinuation}
    \]
    Rechnerisch und geometrisch erhalten wir somit dasselbe Ergebnis.
  \end{minipage}\hfill
  \begin{minipage}[c]{0.44\linewidth}
    \centering
    \CachedTikzFigure{content/resources/chapters/vector_algebra/tikz/vector_subtraction_example_2d.tex}
  \end{minipage}
  \par\smallskip
\end{ExampleBox}
